\chapter{Water, soil, snow and static fields}
\label{cap:suolo}

\section{Precipitation and runoff}
\label{sec:precipitazione}

Precipitation and runoff are fields accumulated over the hour ending at the validity time
(\cref{sec:tempo}). Total precipitation
\campo{tp} (paramId 228, \si{m}) is WRF's resolved precipitation \campo{RAINNC}, which is complete
because the convective precipitation \campo{RAINC} is identically zero (\campo{CU_PHYSICS=0},
\cref{tab:configurazione}). Snowfall \campo{sf} (144, \si{m} of water equivalent) is the sum of
resolved snow and graupel, $\campo{SNOWNC} + \campo{GRAUPELNC}$, that is, all solid precipitation,
as in the ECMWF definition \parencite{paramdb}. Rainfall \campo{tirf} (235015, \si{kg/m^2}) is total
precipitation minus snow and graupel; the microphysics has no hail. The three fields form a closed budget. With
\campo{tirf} converted to metres, $\campo{tp} - \campo{sf} - \campo{tirf}$ is at most
\SI{7.8e-8}{m} at the worst point, on 15 February and 15 July 2024 at 12 UTC.

The three runoff fields, surface (\campo{sro}, paramId 8), subsurface (\campo{ssro}, 9) and total
(\campo{ro}, 205), carry the information of a single field, and the archive publishes all three
\parencite{issue51}. At eight validity times of 2024 (four in \cref{tab:deflusso}), 15 January,
February, March and August at 12 UTC and 15 July at 00, 06, 12 and 18 UTC, \campo{ssro} is
identically zero at four and, at the others, is non-zero at no more than 17 of the \num{2220273}
grid points. At every point of the eight validity times, $|\campo{ro} - \campo{sro} - \campo{ssro}|$
does not exceed \SI{7.5e-9}{m}.

\begin{table}[htbp]
  \centering
  \small
  \caption{Domain means of the runoff fields, accumulated over the hour ending at the validity time,
  in \si{m}, and number of grid points where \campo{ssro} is non-zero.}
  \label{tab:deflusso}
  \begin{tabular}{@{}lSSrS@{}}
    \toprule
    Validity time & {\campo{sro}} & {\campo{ssro}} & Points \campo{ssro} $\neq 0$ & {\campo{ro}} \\
    \midrule
    2024-02-15 12 UTC & 4.819e-6 & 0       & 0  & 4.819e-6 \\
    2024-07-15 12 UTC & 5.543e-6 & 2.2e-9  & 3  & 5.545e-6 \\
    2024-07-15 18 UTC & 1.829e-5 & 1.5e-9  & 2  & 1.829e-5 \\
    2024-08-15 12 UTC & 2.334e-6 & 0       & 0  & 2.334e-6 \\
    \bottomrule
  \end{tabular}
\end{table}

\campo{ssro} is the hourly increment of \campo{UDROFF}, the model's subsurface runoff accumulator.
The accumulator is non-zero at \numrange{462}{8363} points, with maxima of \SIrange{11}{106}{mm}. In the
run started at 18 UTC on 14 July 2024 the same \num{1647} points carry subsurface runoff at its
sixth, twelfth, eighteenth and twenty-fourth hours (00, 06, 12 and 18 UTC on 15 July), only 19 of
them grow between the sixth and the twenty-fourth hour, and in a single hour no more than 17 grow.

Relative to ERA5, with the accumulations of both datasets summed over the hours from 01 UTC to
the validity time, \campo{sro} is \numrange{1.0}{4.6} times, \campo{ssro}
\numrange{0.000}{0.004} times and \campo{ro} \numrange{0.19}{0.85} times the ERA5 value. In
ERA5 the largest term is subsurface runoff, which in the archive is almost everywhere zero
(\cref{tab:deflusso}).

\section{Soil}
\label{sec:suolo}

Volumetric soil water \campo{swvl1-4} (paramId 39--42, \si{m^3/m^3}) and soil temperature
\campo{stl1-4} (139, 170, 183, 236, \si{K}) are given on the four ERA5 layers, 0--7, 7--28, 28--100
and \SIrange{100}{289}{cm}. The RUC scheme instead carries point values at six nodes, at
\SIlist{0;5;20;40;160;300}{cm}, with a linear profile between each node and the next. The four ERA5
layers lie between the first and the last node, because ERA5 starts at \SI{0}{cm}, where the first
node sits, and ends at \SI{289}{cm}, above the node at \SI{300}{cm}. Each layer is therefore the
exact average, over the layer, of the model's linear profile, with no extrapolation at either end.
Each layer is a weighted combination of three consecutive nodes, at 0, 5 and \SI{20}{cm}, at 5, 20
and \SI{40}{cm}, at 20, 40 and \SI{160}{cm} and at 40, 160 and \SI{300}{cm}, with weights that sum
to 1. \campo{swvl} uses the total soil moisture, liquid plus frozen, which is what the ERA5 name
denotes.

The eight fields are defined on the permanent land of the static file, and sea ice has no soil
column, as in ERA5, whose land--sea mask is permanent. Where sea ice forms, WRF reclassifies the
cell as land, assigns it soil type 16 and fills the soil column with water: \campo{swvl} reaches
exactly \num{1.000}, against a porosity of \num{0.435} for that soil type. Those columns are a model
state but not a soil, and the archive masks them. The ice itself stays in the file, in \campo{ci}.

\section{Snow}
\label{sec:neve}

The archive carries the snow water equivalent \campo{sd} (paramId 141, \si{m}), the snow density
\campo{rsn} (33, \si{kg/m^3}), the snow cover \campo{snowc} (260038, \si{\percent}) and the
temperature of the snow layer \campo{tsn} (238, \si{K}). \campo{tsn} is \campo{SOILT1}, the
temperature of the first level of the model's combined snow and soil column, and where snow does
not cover the cell that level is soil. The field is therefore defined only where \campo{SNOWC}
exceeds \num{0.9}, that is, where \campo{snowc} exceeds \SI{90}{\percent}, and is missing elsewhere,
including cells with snow on a smaller fraction. The field is missing everywhere at 7 of the 34
validity times of the verification sample, all of them in summer.

\section{Vegetation, soil type and lakes}
\label{sec:vegetazione}

\paragraph{Cover.} \campo{cvl} (paramId 27) and \campo{cvh} (28) are the sums of the static-file
fractions over the MODIS-IGBP classes of low and high vegetation, the same quantity ERA5 denotes
(\cref{tab:tipi-veg}). High vegetation comprises the IGBP classes with tree cover above
\SI{30}{\percent}, that is, forests (above \SI{60}{\percent}) and woody savannas (between 30 and
\SI{60}{\percent}); savanna, between 10 and \SI{30}{\percent}, is low \parencite{modis2022}.

\paragraph{Vegetation types.} \campo{tvl} (29) and \campo{tvh} (30) are, among the low classes and
among the high ones, the class with the largest fraction in the cell, translated into WMO code table
4.234 following the correspondence in \cref{tab:tipi-veg}. The code table distinguishes evergreen shrubs
(code 16) from deciduous shrubs (17), whereas MODIS does not record shrub phenology. On the two
shrubland classes the code cannot be chosen without an assumption, and \campo{tvl} is missing on
the \num{124672} cells where shrubland dominates, \SI{16.02}{\percent} of the \num{778381} cells
with $\campo{cvl} > 0$. \campo{cvl} still reports the cover in those cells.

\begin{table}[htbp]
  \centering
  \small
  \caption{Translation of the MODIS-IGBP classes into WMO 4.234 codes for \campo{tvl} and
  \campo{tvh}.}
  \label{tab:tipi-veg}
  \begin{tabularx}{\textwidth}{@{}XlX@{}}
    \toprule
    MODIS-IGBP class & Set & 4.234 code \\
    \midrule
    1 evergreen needleleaf forest & high & 3 evergreen needleleaf trees \\
    2 evergreen broadleaf forest & high & 6 evergreen broadleaf trees \\
    3 deciduous needleleaf forest & high & 4 deciduous needleleaf trees \\
    4 deciduous broadleaf forest & high & 5 deciduous broadleaf trees \\
    5 mixed forest & high & 18 mixed forest \\
    8 woody savanna & high & 19 interrupted forest \\
    9 savanna & low & 7 tall grass \\
    10 grassland & low & 2 short grass \\
    11 permanent wetland (absent) & low & 13 bogs and marshes \\
    12 cropland & low & 1 crops, mixed farming \\
    14 cropland and natural vegetation mosaic (absent) & low & 1 crops, mixed farming \\
    18, 19, 20 wooded, mixed and barren tundra & low & 9 tundra \\
    6, 7 closed and open shrubland & low & none: \campo{tvl} missing \\
    13, 15, 16, 17, 21 urban, snow and ice, barren, water, lake & none & none: fraction $1 - \campo{cvl} - \campo{cvh}$ \\
    \bottomrule
  \end{tabularx}
\end{table}

\paragraph{Soil type.} \campo{slt} (paramId 43) does not translate categories. Five of ECMWF's seven
soil types are defined by thresholds on the clay and sand percentages \parencite{ifs41r2p4}, and the
archive applies them to the clay and sand fractions of the static file (\cref{tab:slt}). The
thresholds of the coarse, medium and fine types are those of the FAO \parencite{hwsd2012}; those of
the medium fine and very fine types have no external source. Once the thresholds are applied, type 6
replaces the result where the dominant soil category is organic material. All \num{1304109} permanent-land
cells have a type: \num{1069456} medium, \num{187378} coarse, \num{45450} fine, \num{1498} organic
and 327 very fine. Types 3 (medium fine) and 7 (tropical organic) do not occur. eccodes decodes
\campo{slt} with WMO code table 4.213, which has 11 classes and a different meaning; the values of
the archive must be read with the ECMWF definition.

\begin{table}[htbp]
  \centering
  \small
  \caption{Thresholds for the soil type \campo{slt} (clay and sand in percent).}
  \label{tab:slt}
  \begin{tabular}{@{}ll@{}}
    \toprule
    Condition & ECMWF type \\
    \midrule
    clay $\geq 60$ & 5 very fine \\
    $35 \leq$ clay $< 60$ & 4 fine \\
    clay $< 35$ and sand $\leq 15$ & 3 medium fine \\
    $18 \leq$ clay $< 35$ and sand $> 15$, or clay $< 18$ and $15 <$ sand $\leq 65$ & 2 medium \\
    clay $< 18$ and sand $> 65$ & 1 coarse \\
    dominant soil category of organic material & 6 organic \\
    \bottomrule
  \end{tabular}
\end{table}

\paragraph{Lakes.} The lake cover \campo{cl} (paramId 26) comes from the fraction of the lake class
in the static file, and it is non-zero on \num{32000} cells. It is not redundant with the hourly
land--sea mask \campo{lsm}: \num{19309} cells that the mask counts as land contain a lake smaller
than the cell. The lake depth \campo{dl} (228007, \si{m}) is the \campo{LAKE_DEPTH} field of the
static file, and it is not a model state, because the lake physics is switched off
(\campo{SF_LAKE_PHYSICS=0}). \campo{dl} is defined on the same \num{32000} cells as \campo{cl} and
is missing elsewhere; on \num{22991} of them it is \SI{10}{m}, the WPS default value. In the static
file \campo{LAKE_DEPTH} takes this default value on \SI{99.56}{\percent} of the domain.

\section{Albedo}
\label{sec:albedo}

The background albedo \campo{al} (paramId 174) is a table value per dominant category. The
simulation does not use monthly satellite albedos, and the value comes from the winter and summer
columns of the land-use table: the summer column applies from April to October, the winter column
from November to March. Over the whole domain the field takes 8 to 11 distinct values, and at the
switch between the two columns it changes on \num{1171937} cells, \SI{89.9}{\percent} of the
permanent-land cells. Each change is one row of the table: for example from \num{0.20} to
\num{0.17} on cropland on \num{497649} cells and from \num{0.23} to \num{0.25} on barren soil on
\num{369219}. Over water it is \num{0.08}, over sea ice \num{0.65}. Within a day the value over land
does not change; it changes only on sea-ice cells, which follow the hourly mask
(\cref{sec:statici-orari}). \campo{al} is neither an observed albedo nor an evolving one, and a
difference between two dates on either side of the column switch is not a physical signal.

The effective albedo of the simulation is \campo{fal} (paramId 243), the ratio of reflected to
incident solar radiation at the surface. It is defined only where the incident radiation exceeds
\SI{50}{W/m^2}, and is missing elsewhere.

\section{Orography}
\label{sec:orografia}

The orography is the surface geopotential \campo{z} (paramId 129). The slope \campo{slor} (163) is
the magnitude of the gradient of the terrain height $h$ on the \SI{3}{km} grid, $|\nabla h|$, whereas
paramId 163 denotes the slope of the sub-grid orography. eccodes shows the parameter name,
\emph{Slope of sub-gridscale orography}, and nothing in the message states that the field is the
resolved slope. The standard deviation of orography \campo{sdor} (160), computed by WPS, is the
ERA5 quantity, evaluated over a \SI{3}{km} cell instead of a \SI{31}{km} one. The ratio of medians
against ERA5 is \num{0.244}, equal to $(3/31)^{0.605}$, and the difference comes from the
resolution.
